\documentclass[12pt]{article}
\usepackage[utf8]{inputenc}
\usepackage[T1]{fontenc}
\usepackage[margin=1in]{geometry}
\usepackage{setspace}
\setlength{\parindent}{2em}
\setlength{\parskip}{1.25em}
\renewcommand{\baselinestretch}{1.0}
\makeatletter
\newcommand{\@subtitle}{}
\newcommand{\subtitle}[1]{\renewcommand{\@subtitle}{#1}}
\renewcommand{\@maketitle}{%
  \newpage
  \null
  \vskip 1em%
  \begin{center}%
    {\LARGE \@title \par}%
    \ifx\@subtitle\@empty\else
      \vskip 0.5em%
      {\large \@subtitle \par}%
    \fi
  \end{center}%
  \par
  \vskip 1em}
\makeatother
\usepackage{xcolor}
\usepackage[colorlinks=true,urlcolor=blue!55!black,linkcolor=black]{hyperref}

\title{\textit{Cruzan v. Director, Missouri Department of Health}, 497 U.S. 261 (1990)}
\subtitle{Case Brief}
\author{}
\date{}

\begin{document}
\maketitle

\section{Memory Jogger}

In \href{https://www.courtlistener.com/opinion/112478/cruzan-ex-rel-cruzan-v-director-missouri-department-of-health/}{\textit{Cruzan v. Director, Missouri Department of Health}}, the Supreme Court assumed that a competent person has a Fourteenth Amendment liberty interest in refusing lifesaving nutrition and hydration, but held that a State may constitutionally require ``clear and convincing evidence'' of an incompetent patient's wishes before permitting family members to withdraw life support.

\section{Detailed Case Facts}

On the night of January 11, 1983, 25-year-old Nancy Beth Cruzan lost control of her car on Elm Road in Jasper County, Missouri. Her car overturned, throwing her face-down into a ditch without detectable respiratory or cardiac function (\href{https://supreme.justia.com/cases/federal/us/497/261/}{497 U.S. at 266}). Paramedics restored her heartbeat and breathing at the scene, but she had suffered anoxia (oxygen deprivation) for an estimated 12 to 14 minutes---well beyond the 6 minutes after which permanent brain damage typically occurs. Following roughly three weeks in a coma, she progressed to a persistent vegetative state (PVS), a condition in which a patient retains motor reflexes and spontaneous respiration and circulation but exhibits no significant cognitive function and has irreversible, permanent, progressive cerebral cortical atrophy (\href{https://supreme.justia.com/cases/federal/us/497/261/}{497 U.S. at 266 \& n.1}).

With the consent of her then-husband, surgeons implanted a gastrostomy feeding and hydration tube to sustain her. Rehabilitative efforts failed; Nancy remained a spastic quadriplegic with no ability to swallow and virtually no chance of regaining cognitive faculties, though medical experts testified she could be kept biologically alive for another 30 years. She lay in a Missouri state hospital, with the State of Missouri bearing the cost of her care (\href{https://supreme.justia.com/cases/federal/us/497/261/}{497 U.S. at 266--267 \& n.1}).

Once it became clear Nancy would never recover, her parents and co-guardians, Lester and Joyce Cruzan, asked hospital employees to terminate her artificial nutrition and hydration. Because all parties agreed removal of the feeding tube would cause her death, the hospital refused without court approval. The litigation centered on the Due Process Clause of the Fourteenth Amendment (``nor shall any State deprive any person of life, liberty, or property, without due process of law''), Missouri's Living Will statute, Mo. Rev. Stat. \S\ 459.010 \textit{et seq.} (1986), and Missouri's statutory definition of death, Mo. Rev. Stat. \S\ 194.005 (1986) (\href{https://supreme.justia.com/cases/federal/us/497/261/}{497 U.S. at 266 n.1, 268, 278}).

Historically and politically, \textit{Cruzan} was the Supreme Court's first ``right to die'' case, arriving fourteen years after the New Jersey Supreme Court's landmark decision in \textit{In re Quinlan}, 70 N.J. 10, 355 A.2d 647 (1976). Advances in resuscitation and enteral feeding had created a new population of roughly 10,000 Americans maintained in persistent vegetative states, sparking intense legal, bioethical, and religious debate alongside the Court's broader battle over substantive due process following \textit{Roe v. Wade} (1973) and \textit{Bowers v. Hardwick} (1986). The case catalyzed widespread public demand for living wills and health care proxies and helped prompt Congress's enactment of the Patient Self-Determination Act of 1990 (\href{https://scholarship.law.nd.edu/cgi/viewcontent.cgi?article=1280&context=jleg}{Notre Dame J. Legis.}).

\section{Procedural History}

\begin{itemize}
  \item \textbf{Missouri Trial Court (Circuit Court of Jasper County, Probate Division):} The trial court authorized Nancy's parents to terminate artificial nutrition and hydration, ruling that a person in Nancy's condition had a fundamental right under both the Missouri and Federal Constitutions to refuse or direct the withdrawal of ``death prolonging procedures.'' Relying primarily on testimony from a former housemate (Athena Comer) that Nancy had stated at age 25 in a ``somewhat serious conversation'' that if sick or injured she would not wish to continue living unless she could live ``at least halfway normally,'' the trial court found that the evidence ``suggest[ed]'' she would not wish to continue artificial feeding (\href{https://supreme.justia.com/cases/federal/us/497/261/}{497 U.S. at 268, 285}).
  \item \textbf{Supreme Court of Missouri:} On appeal by the State of Missouri and Nancy's court-appointed guardian ad litem, the Missouri Supreme Court reversed en banc by a 4--3 vote in \textit{Cruzan v. Harmon}, 760 S.W.2d 408 (Mo. 1988). While recognizing the common-law doctrine of informed consent, the court doubted its applicability, declined to find a broad right of privacy in either the Missouri or Federal Constitutions, read Missouri's Living Will statute as expressing a strong state policy favoring preservation of life, dismissed Nancy's statements to her housemate as ``unreliable,'' and held that no surrogate could terminate treatment for an incompetent patient absent living-will formalities or ``clear and convincing, inherently reliable evidence'' of the patient's wishes (\href{https://supreme.justia.com/cases/federal/us/497/261/}{497 U.S. at 268--269}).
  \item \textbf{U.S. Supreme Court:} Granted certiorari, 492 U.S. 917 (1989); argued December 6, 1989; affirmed (5--4) on June 25, 1990 (\href{https://tile.loc.gov/storage-services/service/ll/usrep/usrep497/usrep497261/usrep497261.pdf}{497 U.S. at 261, 287}). (In subsequent state proceedings after Missouri withdrew from the case in September 1990, Nancy's parents presented new testimony from former coworkers; Jasper County Probate Judge Charles Teel found clear and convincing evidence of Nancy's wishes, the feeding tube was removed on December 14, 1990, and Nancy died on December 26, 1990 (\href{https://scholarship.law.nd.edu/cgi/viewcontent.cgi?article=1280&context=jleg}{Notre Dame J. Legis.}).)
\end{itemize}

\section{Judicial Votes}

\begin{itemize}
  \item \textbf{Majority (5--4):} Chief Justice William H. Rehnquist delivered the opinion of the Court, joined by Justices Byron R. White, Sandra Day O'Connor, Antonin Scalia, and Anthony M. Kennedy (\href{https://supreme.justia.com/cases/federal/us/497/261/}{497 U.S. at 263}).
  \item \textbf{Concurrences:}
    \begin{itemize}
      \item Justice Sandra Day O'Connor filed a concurring opinion (joining the Court's opinion).
      \item Justice Antonin Scalia filed a concurring opinion (joining the Court's opinion).
    \end{itemize}
  \item \textbf{Dissents:}
    \begin{itemize}
      \item Justice William J. Brennan, Jr. filed a dissenting opinion, joined by Justices Thurgood Marshall and Harry A. Blackmun.
      \item Justice John Paul Stevens filed a dissenting opinion.
    \end{itemize}
\end{itemize}

\section{Holding}

\textbf{Judgment:} Affirmed the judgment of the Supreme Court of Missouri (\href{https://supreme.justia.com/cases/federal/us/497/261/}{497 U.S. at 287}).

\textbf{Rules Announced:}
\begin{enumerate}
  \item \textbf{Due Process Liberty Interest in Refusing Treatment:} A competent person has a constitutionally protected liberty interest under the Fourteenth Amendment Due Process Clause in refusing unwanted medical treatment, to be balanced against relevant state interests (\href{https://supreme.justia.com/cases/federal/us/497/261/}{497 U.S. at 278--279}). For purposes of this case, the Court assumed that a competent person has a constitutionally protected right to refuse lifesaving artificial hydration and nutrition.
  \item \textbf{Permissibility of Clear and Convincing Evidence Standard:} The Federal Constitution does not forbid a State from requiring ``clear and convincing evidence'' of an incompetent patient's previously expressed wishes before permitting a surrogate to withdraw life-sustaining treatment. Given the State's unqualified interest in preserving human life, guarding against surrogate abuse or non-adversarial proceedings, and the asymmetric irrevocability of an erroneous decision to terminate life support, a State may place the heightened risk of error on those seeking withdrawal (\href{https://supreme.justia.com/cases/federal/us/497/261/}{497 U.S. at 280--285}).
  \item \textbf{No Constitutional Mandate for Family Substituted Judgment:} The Due Process Clause does not require a State to accept the ``substituted judgment'' of close family members absent substantial proof that their views reflect the patient's own wishes (\href{https://supreme.justia.com/cases/federal/us/497/261/}{497 U.S. at 285--287}).
  \item \textbf{Application to Record:} The Missouri Supreme Court did not commit constitutional error in concluding that Nancy Cruzan's informal statements about not wanting to live as a ``vegetable'' did not constitute clear and convincing proof of a desire to withdraw artificial nutrition and hydration (\href{https://supreme.justia.com/cases/federal/us/497/261/}{497 U.S. at 285}).
\end{enumerate}

\textbf{Holding vs. Dicta:} The Court's actual holding is narrow and procedural: the Due Process Clause permits a State to impose a clear-and-convincing-evidence requirement on surrogate requests to withdraw life-sustaining treatment and to decline uncorroborated family substituted judgment. Crucially, the majority only \textit{assumed}---rather than formally held---that a competent person has a constitutional right to refuse lifesaving artificial nutrition and hydration (\href{https://supreme.justia.com/cases/federal/us/497/261/}{497 U.S. at 279}), although five Justices (Justice O'Connor concurring plus the four dissenters) explicitly affirmed that right, and the Court later treated that right as grounded in common-law bodily integrity in \href{https://supreme.justia.com/cases/federal/us/521/702/}{\textit{Washington v. Glucksberg}, 521 U.S. 702 (1997)}. In footnote 12, the Court also expressly reserved whether a State must honor the decision of a surrogate whom the patient formally designated while competent.

\section{Analysis of Opinions}

\subsection{Majority Opinion (Chief Justice Rehnquist)}

Chief Justice Rehnquist began by surveying the common-law doctrine of informed consent (\textit{Union Pacific R. Co. v. Botsford}; \textit{Schloendorff v. Society of New York Hospital}) and the diverse body of post-\textit{Quinlan} state court decisions (\href{https://supreme.justia.com/cases/federal/us/497/261/}{497 U.S. at 269--278}). Emphasizing that federal courts lack the common-law and state-statutory tools available to state courts, Rehnquist framed the issue narrowly: whether the Federal Constitution forbids Missouri's rule of decision.

Citing \textit{Jacobson v. Massachusetts}, \textit{Washington v. Harper}, \textit{Vitek v. Jones}, and \textit{Parham v. J.R.}, Rehnquist recognized that a competent person has a Fourteenth Amendment ``liberty interest'' (expressly rejecting a generalized ``right of privacy'' framework in footnote 7) in refusing unwanted medical treatment, which must be balanced against relevant state interests under \textit{Youngberg v. Romeo}. Without resolving whether artificial feeding is constitutionally identical to other medical interventions, the Court assumed for purposes of the case that a competent adult has a right to refuse lifesaving hydration and nutrition (\href{https://supreme.justia.com/cases/federal/us/497/261/}{497 U.S. at 278--279}).

An incompetent patient, however, cannot make an informed, voluntary choice; any refusal must be exercised by a surrogate. Rehnquist held that Missouri could legitimately require clear and convincing evidence of the patient's wishes to advance several weighty state interests: (1) an unqualified interest in preserving human life, without making ``quality of life'' judgments; (2) safeguarding the personal element of life-and-death decisions; (3) protecting vulnerable patients against potential family abuse; and (4) compensating for the fact that surrogate-termination hearings may not be truly adversarial (noting in footnote 9 that even Nancy's guardian ad litem favored termination). Drawing on \textit{Addington v. Texas} and \textit{Santosky v. Kramer}, Rehnquist reasoned that standards of proof allocate the risk of error: an erroneous decision \textit{not} to terminate preserves the status quo (leaving room for medical advances, new evidence, or legal changes), whereas an erroneous decision \textit{to} terminate is irrevocable (\href{https://supreme.justia.com/cases/federal/us/497/261/}{497 U.S. at 280--284}). Analogizing to statutes of wills and statutes of frauds, Rehnquist observed that general evidentiary rules need not work faultlessly. Finally, he rejected petitioners' alternative claim for family ``substituted judgment,'' explaining that cases permitting States to defer to families (\textit{Michael H. v. Gerald D.}; \textit{Parham}) do not \textit{compel} States to do so (\href{https://supreme.justia.com/cases/federal/us/497/261/}{497 U.S. at 285--287}).

\subsection{Concurring Opinion (Justice O'Connor)}

Justice O'Connor supplied the pivotal fifth vote while writing separately to make two major points (\href{https://supreme.justia.com/cases/federal/us/497/261/}{497 U.S. at 287--292}). First, she refused to leave the underlying right merely ``assumed'': because artificial feeding via nasogastric or surgically implanted gastrostomy tubes involves physical intrusion and restraint, ``the liberty guaranteed by the Due Process Clause must protect, if it protects anything, an individual's deeply personal decision to reject medical treatment, including the artificial delivery of food and water.'' Second, highlighting that few Americans execute detailed written living wills, O'Connor emphasized that the Court's opinion did not decide whether a State has a constitutional duty to honor the decision of a patient's duly appointed health care proxy or durable power of attorney---signaling that such a duty ``may well be constitutionally required'' while entrusting procedural experimentation to the ``laboratory'' of the States.

\subsection{Concurring Opinion (Justice Scalia)}

Justice Scalia joined Rehnquist's opinion but wrote separately to urge complete judicial abstention: ``the federal courts have no business in this field'' (\href{https://supreme.justia.com/cases/federal/us/497/261/}{497 U.S. at 292--301}). Applying a strict historical-tradition test for substantive due process (\textit{Michael H.}; \textit{Bowers}), Scalia argued that American and English common law had always empowered the State to prevent suicide, by force if necessary. He rejected all three distinctions offered by petitioners: (1) common-law suicide prohibitions made no exception for terminal illness or suffering; (2) the action/inaction distinction was legally specious (starving oneself to death or refusing to come in from the cold is as much suicide as using a gun, and drawing a line between ``ordinary'' and ``extraordinary'' care is a legislative policy judgment, not a constitutional rule); and (3) bodily integrity cannot override the State's traditional authority to prevent self-destruction. Warning against repeating the judicial confusion of abortion jurisprudence, Scalia concluded that the Equal Protection Clause---requiring majorities to live under the same rules they impose on others---provides the real check against oppressive end-of-life laws.

\subsection{Dissenting Opinion (Justice Brennan, joined by Justices Marshall and Blackmun)}

Justice Brennan argued that freedom from unwanted medical treatment---including surgically implanted feeding tubes, which the medical profession, Medicare, and the FDA all treat as medical interventions---is a \textit{fundamental} right deeply rooted in common-law bodily integrity (\href{https://supreme.justia.com/cases/federal/us/497/261/}{497 U.S. at 301--330}). Because a fundamental right was burdened, Missouri's evidentiary barrier triggered strict scrutiny under \textit{Zablocki v. Redhail}. Brennan maintained that an incompetent patient retains that fundamental right, and that the State has no legitimate general interest in a person's life ``completely abstracted from the interest of the person living that life.'' Consequently, Missouri's only legitimate interest before Nancy's wishes were determined was \textit{accuracy}.

Brennan dismantled the majority's defense of Missouri's asymmetrical clear-and-convincing rule: unlike \textit{Santosky} and \textit{Addington}, where heightened proof shielded an individual's fundamental right against the State, Missouri deployed clear and convincing evidence as an \textit{obstacle} to a fundamental right. Moreover, error is irrevocable in \textit{both} directions: erroneously keeping an irreversibly vegetative patient tethered to machines degrades her bodily dignity, protracts her family's anguish, and distorts her memory. By demanding formal written directives that fewer than 15--23\% of Americans execute and disregarding serious conversations with family and friends, Missouri skewed factfinding away from accuracy and unlawfully appropriated the patient's choice rather than entrusting it to the family who knew her best.

\subsection{Dissenting Opinion (Justice Stevens)}

Justice Stevens wrote separately to focus on Nancy Cruzan's objective ``best interests'' rather than merely her prior statements (\href{https://supreme.justia.com/cases/federal/us/497/261/}{497 U.S. at 330--357}). Because the trial judge, Nancy's loving parents, and the independent guardian ad litem all agreed that discontinuing tube feeding served Nancy's best interests, Stevens argued that Missouri's rule improperly treated the absence of a living will as a ``waiver'' of constitutional liberty---leaving children, accident victims, and the elderly unprotected. More fundamentally, Stevens contended that ``[c]hoices about death touch the core of liberty'' and that Missouri was not preserving life but imposing an aberrant, ``sectarian definition of life'' that reduced human existence to mere biological persistence severed from the living person's own interests. Abstracting Nancy's biological body from her personhood to serve a symbolic state teaching function violated the Due Process Clause.

\section{Examples: Future Applications}

\begin{enumerate}
  \item \textbf{Same Side (Constitutional State Restriction):} A State enacts a statute prohibiting surrogates from withdrawing mechanical ventilation from an irreversibly comatose patient unless the patient signed a notarized advance directive or made specific, corroborated oral statements proven by clear and convincing evidence. \textit{Why:} Directly governed by \textit{Cruzan}, the State may allocate the risk of error toward preserving life when an incompetent patient's wishes are uncertain, even if close relatives uniformly believe the patient would refuse ventilation.
  \item \textbf{Same Side (Constitutional State Restriction):} A State prohibits physicians from prescribing lethal medication to a competent, terminally ill adult who wishes to end her life (physician-assisted suicide). \textit{Why:} Consistent with \textit{Cruzan}'s narrow grounding in common-law bodily integrity against unwanted medical intrusion---rather than a broad constitutional ``right to die''---the Supreme Court later upheld such bans in \href{https://supreme.justia.com/cases/federal/us/521/702/}{\textit{Washington v. Glucksberg}, 521 U.S. 702 (1997)}, distinguishing refusal of unwanted treatment from affirmative assistance in suicide.
  \item \textbf{Opposite Side (Unconstitutional State Action):} A state hospital forcibly inserts a feeding tube into a fully competent, mentally sound 70-year-old cancer patient who explicitly and repeatedly refuses artificial nutrition after being informed it will prolong her life by several months. \textit{Why:} Five Justices in \textit{Cruzan} (O'Connor plus the four dissenters) explicitly recognized, and the majority assumed, a competent adult's Due Process liberty interest in refusing unwanted medical treatment including artificial feeding; absent third-party harms or public-health threats like contagious disease, the State's abstract interest in preserving life cannot override a competent patient's contemporaneous refusal.
  \item \textbf{Opposite Side (Unconstitutional State Action):} A State passes a law declaring that even when a patient executes a valid, unambiguous written living will specifically directing that artificial nutrition and hydration be withheld in the event of a persistent vegetative state, state hospitals must ignore the directive and maintain tube feeding indefinitely. \textit{Why:} Unlike Missouri's procedural safeguard designed to verify the patient's actual wishes, a categorical refusal to honor clear, formal, written instructions directly extinguishes the substantive liberty interest recognized in \textit{Cruzan} without any evidentiary justification.
  \item \textbf{Fence-Sitter (Borderline):} While competent, a patient executes a valid state durable power of attorney for health care naming her sister as proxy with full authority to make end-of-life decisions, but leaves no written instructions mentioning feeding tubes; after the patient enters a persistent vegetative state, the State refuses to let the sister withdraw a gastrostomy tube because there is no clear and convincing evidence of the patient's specific treatment preference. \textit{Why:} Footnote 12 of Rehnquist's majority opinion explicitly left open whether a State must defer to a duly appointed health care surrogate, while Justice O'Connor's pivotal concurrence strongly suggested that honoring a formally designated proxy ``may well be constitutionally required'' to make the liberty interest meaningful.
\end{enumerate}

\section{Critique}

\textit{Cruzan} was widely praised for stimulating advance-directive legislation and avoiding an abrupt national constitutional rule, yet both legal scholars and clinicians have identified deep doctrinal and practical flaws (\href{https://www.acpjournals.org/doi/10.7326/0003-4819-114-10-895}{Annals of Internal Medicine}).

\textbf{Logical Weaknesses and Competing Premises:} At the core of \textit{Cruzan} is an unresolved tension between procedural neutrality and substantive value choice. Chief Justice Rehnquist characterized Missouri's clear-and-convincing standard as a neutral ``procedural safeguard'' that preserves the ``status quo.'' Critics, echoing Justice Brennan and bioethicist Nancy Rhoden, point out that keeping a permanently unconscious patient biologically alive via surgically implanted tubes is not a natural ``status quo''---it is an affirmative, ongoing technological intervention. Furthermore, Rehnquist's asymmetry argument (that erroneous continuation of life support is reversible while erroneous termination is final) collapses for a patient with irreversible cortical destruction: years of unwanted bodily maintenance and family suffering cannot be undone. Finally, by invoking \textit{Addington} and \textit{Santosky}---cases where heightened proof protected individuals \textit{against} state coercion---to justify a burden that empowers the State to override family consensus, the majority inverted the protective function of Due Process proof standards.

\textbf{Originalist and Textualist Perspective:} From an originalist and textualist perspective, best articulated by Justice Scalia's concurrence, the majority opinion was insufficiently disciplined because it accepted (or assumed) an unenumerated substantive due process liberty interest in refusing lifesaving nutrition without grounding it in 1868 constitutional text or history. Originalists contend that ``due process of law'' guarantees fair procedures, not substantive bioethical rights; that common-law battery law protected against unauthorized medical touching by private physicians rather than stripping States of authority to prevent self-starvation; and that courts lack any objective constitutional metric to distinguish ``ordinary'' sustenance from ``extraordinary'' medical treatment (\href{https://journals.library.columbia.edu/index.php/bioethics/article/view/14627}{Columbia Bioethics}; \href{https://texaslawreview.org/wp-content/uploads/2025/02/Faulkner.Printer.pdf}{Texas L. Rev.}). By leaving substantive end-of-life boundaries half-open---particularly through Justice O'Connor's concurrence---the Court invited protracted federal litigation over questions better resolved by state legislatures.

\textbf{Progressive and Patient-Autonomy Perspective:} Progressive, bioethical, and feminist scholars criticize \textit{Cruzan} from the opposite vantage point: for bowing to conservative anxieties over substantive due process and abandoning vulnerable patients who lack the legal sophistication or foresight to draft formal living wills. Critics note that the Court downgraded a fundamental right of bodily autonomy into a weaker ``liberty interest'' subject to amorphous balancing, effectively privileging wealthy or legally advised citizens who execute advance directives while condemning the overwhelming majority---especially young accident victims like Nancy Cruzan, children, and the never-competent---to indefinite technological maintenance whenever a State asserts an abstract ``sanctity of life'' policy. Medical commentators also warned that \textit{Cruzan}'s rigid evidentiary standard could undermine trust between families and physicians, encourage defensive medicine, and perversely deter doctors from initiating potentially lifesaving emergency interventions out of fear they could never legally be withdrawn (\href{https://www.acpjournals.org/doi/10.7326/0003-4819-114-10-895}{Annals of Internal Medicine}).

\section{Key Quotations}

\begin{quote}
An erroneous decision not to terminate results in a maintenance of the status quo; the possibility of subsequent developments such as advancements in medical science, the discovery of new evidence regarding the patient's intent, changes in the law, or simply the unexpected death of the patient despite the administration of life-sustaining treatment, at least create the potential that a wrong decision will eventually be corrected or its impact mitigated. An erroneous decision to withdraw life-sustaining treatment, however, is not susceptible of correction.
--- Chief Justice William H. Rehnquist, Opinion of the Court (\href{https://supreme.justia.com/cases/federal/us/497/261/}{497 U.S. at 283})
\end{quote}

\begin{quote}
Requiring a competent adult to endure such procedures against her will burdens the patient's liberty, dignity, and freedom to determine the course of her own treatment. Accordingly, the liberty guaranteed by the Due Process Clause must protect, if it protects anything, an individual's deeply personal decision to reject medical treatment, including the artificial delivery of food and water.
--- Justice Sandra Day O'Connor, Concurring (\href{https://supreme.justia.com/cases/federal/us/497/261/}{497 U.S. at 289})
\end{quote}

\begin{quote}
While I agree with the Court's analysis today, and therefore join in its opinion, I would have preferred that we announce, clearly and promptly, that the federal courts have no business in this field; that American law has always accorded the State the power to prevent, by force if necessary, suicide---including suicide by refusing to take appropriate measures necessary to preserve one's life; that the point at which life becomes ``worthless,'' and the point at which the means necessary to preserve it become ``extraordinary'' or ``inappropriate,'' are neither set forth in the Constitution nor known to the nine Justices of this Court any better than they are known to nine people picked at random from the Kansas City telephone directory; and hence, that even when it is demonstrated by clear and convincing evidence that a patient no longer wishes certain measures to be taken to preserve her life, it is up to the citizens of Missouri to decide, through their elected representatives, whether that wish will be honored.
--- Justice Antonin Scalia, Concurring (\href{https://supreme.justia.com/cases/federal/us/497/261/}{497 U.S. at 293})
\end{quote}

\begin{quote}
But the State has no legitimate general interest in someone's life, completely abstracted from the interest of the person living that life, that could outweigh the person's choice to avoid medical treatment.
--- Justice William J. Brennan, Jr., Dissenting (\href{https://supreme.justia.com/cases/federal/us/497/261/}{497 U.S. at 313})
\end{quote}

\begin{quote}
Nancy Cruzan is obviously ``alive'' in a physiological sense. But for patients like Nancy Cruzan, who have no consciousness and no chance of recovery, there is a serious question as to whether the mere persistence of their bodies is ``life'' as that word is commonly understood, or as it is used in both the Constitution and the Declaration of Independence. The State's unflagging determination to perpetuate Nancy Cruzan's physical existence is comprehensible only as an effort to define life's meaning, not as an attempt to preserve its sanctity.
--- Justice John Paul Stevens, Dissenting (\href{https://supreme.justia.com/cases/federal/us/497/261/}{497 U.S. at 345})
\end{quote}

\end{document}
